\scanpage{32}
\mspage{6)}

il n'est pas possible d'y introduire une métrique riemannienne canonique, même modulo $\mathrm c^{\mathrm{te}}$. Considérons cependant la donnée d'une structure supplémentaire du type
\[
H\subset\mathrm{Aut}_{\mathrm{conf}}(X)=G
\]
où $H$ est un sous-groupe \uline{compact} de $\mathrm{Aut}_{\unsure{G|0}\,\mathrm{conf}}(X)$. Condition \uline{nécessaire} pour que $H$, supposé fermé, invarie une métrique euclidienne sphérique comp. avec la str. conforme). Soit $H^+=H\cap\mathrm{Aut}^+(X)$ (ss-groupe des autom qui conservent l'orientation). OPS $X=\mathbb P(V)$, $V$ vectoriel complexe de rg $2$, \deleted{\unsure{alors|étant}} soit $\overset{\deleted{\unsure{\sim}}}{\widetilde H}$ l'image inverse de $H^+$ dans $\unsure{\mathrm{SL}_{\mathbb C}(2,V)|\mathrm{SL}(2,V)|\mathrm{SL}(V)}$, c'est donc un sous-groupe cpt, qui invarie nécessairement une forme hermitienne déf. pos. sur $V$, donc la métrique sphérique corr. sur $X$ \inserted{(normalisée par la condition $\int_X dx=1$)}. Dans le cas où \deleted{cette métrique} \inserted{la métrique} hermitienne invarie est unique mod $\mathrm c^{\mathrm{te}}$, on trouve donc une métrique \inserted{sphérique normalisée} canonique sur $X$, définie par la \inserted{seule} donnée de $H^+$. Cette métrique sera $=$ invariante par $\mathrm{Norm}_G(H^+)\deleted{\unsure{\subset\mathrm{Aut}}}$, et a fortiori par $H$, et est encore caractérisée par cette condition.

% archive sheet number: 34
